%%%%%%%%%%%%%%%%%%%%%%%%%%%%%%%%%%%%%%%%%%%%%%%%%%%%%%%%%%%%%%%%%%%%%%%%%%%%%%%
%%  Public-Data Reproduction of GIGA-Lens Modeling for the DESI Strong Lens
%%  Foundry Demonstration System DESI-165.4754-06.0423
%%
%%  Two-column AASTeX 7.0.1 manuscript, in preparation for ApJ (2026).
%%%%%%%%%%%%%%%%%%%%%%%%%%%%%%%%%%%%%%%%%%%%%%%%%%%%%%%%%%%%%%%%%%%%%%%%%%%%%%%
\documentclass[twocolumn,linenumbers]{aastex701}

\usepackage{amsmath}
\usepackage{amssymb}
\usepackage{booktabs}
\usepackage{graphicx}
\usepackage{hyperref}

% Convenience macros
\newcommand{\thetaE}{\ensuremath{\theta_{\mathrm{E}}}}
\newcommand{\gammaEPL}{\ensuremath{\gamma_{\mathrm{EPL}}}}
\newcommand{\gext}{\ensuremath{\gamma_{\mathrm{ext}}}}
\newcommand{\gextone}{\ensuremath{\gamma_{\mathrm{ext},1}}}
\newcommand{\gexttwo}{\ensuremath{\gamma_{\mathrm{ext},2}}}
\newcommand{\eone}{\ensuremath{e_{1}}}
\newcommand{\etwo}{\ensuremath{e_{2}}}
\newcommand{\foundry}{DESI--165.4754$-$06.0423}
\newcommand{\gigalens}{\textsc{GIGA-Lens}}
\newcommand{\lenstronomy}{\textsc{lenstronomy}}
\newcommand{\photutils}{\textsc{photutils}}

\received{2026 May}
\revised{}
\accepted{In preparation; to be submitted to ApJ}

\shorttitle{Reproducing DESI Foundry I on \foundry}
\shortauthors{Benson \& Huang}

\begin{document}

\title{Public-Data Reproduction of GIGA-Lens Modeling for the DESI Strong
Lens Foundry Demonstration System DESI--165.4754$-$06.0423}

\correspondingauthor{Greg Benson}

\author{Greg Benson}
\affiliation{Department of Computer Science,
University of San Francisco,
San Francisco, CA 94117, USA}
\email{benson@cs.usfca.edu}

\author{Xiaosheng Huang}
\affiliation{Department of Physics and Astronomy,
University of San Francisco,
San Francisco, CA 94117, USA}
\affiliation{Physics Division,
Lawrence Berkeley National Laboratory,
Berkeley, CA 94720, USA}
% TODO confirm preferred contact email for co-author before submission
\email{xhuang@usfca.edu}

\begin{abstract}
We present an end-to-end, public-data reproduction of the GIGA-Lens
modeling demonstration for the DESI Strong Lens Foundry I system
\foundry, originally reported by \citet{huang2025foundry1}.
Starting from the same {\it Hubble Space Telescope}\ WFC3/F140W exposure
(GO--15867; MAST DOI 10.17909/hx0v-9260) and using only open-source
infrastructure --- \gigalens, \lenstronomy, \photutils, JAX, and
\textsc{TensorFlow Probability} --- we rebuild the data products,
construct an empirical PSF from in-field stars, and fit an EPL$+$external
shear mass model together with four S\'{e}rsic lens-light components,
a S\'{e}rsic source, and a shapelet basis of order $n_{\max}=6$
($74$ free non-linear parameters). Hamiltonian Monte Carlo did not
JIT-compile at this model size under JAX\,0.6.2; switching to the
No-U-Turn Sampler (NUTS) with a static \texttt{max\_tree\_depth} resolves
the compile-time blocker (NUTS compiles in $\sim$5\,min and runs 700
chain steps in 2--3\,min), but the chain we obtain is a diagnostic,
not a usable posterior: the v10 SVI covariance is rank-deficient
(rank 53/74, minimum eigenvalue $-4.4\times10^{-13}$) and silently
yields a NaN Cholesky in JAX, and once jittered the dual-averaging
adapter consistently drives the step size to $\sim$2--$3\times10^{-6}$
so the chain barely moves. We therefore report the v10 Gaussian
stochastic variational posterior (initialized from the best paper-mode
chain of a 200-chain multi-start MAP search) as the headline result, and
present a v11 NUTS run started from the v7 paper-mode MAP point as a
diagnostic that confirms posterior multimodality and locates a third
distinct local mode. Our 10\,000-sample variational posterior recovers
all six sign quadrants of the mass model, matches \thetaE\ to 3.0\%,
$e_{1}$ to 2.5\%, and the external shear position angle to within
$1^{\circ}$ of the published HMC result; the v11f NUTS mode reproduces
\thetaE\ to 1.6\%, the closest of any version to date. The remaining
gap --- a $12^{\circ}$ rotation of the lens ellipticity vector,
$\sim30\%$ smaller $|\etwo|$ and $|\gext|$, and a $57\%$ steeper
deprojected slope $\gammaEPL=2.15$ versus the published $1.37$ --- is
shown to be coherent with a single missing constraint, plausibly a
mass--light position-angle tie or higher-order multipole correction.
We document upstream issues in \gigalens\,2.0/JAX\,0.6.2/TFP\,0.25.0
suitable for filing as GitHub issues, and provide the full software
stack table to support future Foundry-series reproductions.
\end{abstract}

\keywords{Strong gravitational lensing (1643) ---
Computational methods (1965) --- Open source software (1866) ---
Reproducibility}

%%%%%%%%%%%%%%%%%%%%%%%%%%%%%%%%%%%%%%%%%%%%%%%%%%%%%%%%%%%%%%%%%%%%%%%%%%%%%%%
\section{Introduction}
\label{sec:intro}
%%%%%%%%%%%%%%%%%%%%%%%%%%%%%%%%%%%%%%%%%%%%%%%%%%%%%%%%%%%%%%%%%%%%%%%%%%%%%%%

Strong gravitational lensing by galaxies is one of the rare astrophysical
phenomena that maps directly to a small number of free physical
parameters --- the projected mass profile of the deflector, the
source position and morphology, and (in some configurations) the
cosmological distance ratio $D_{ls}/D_{s}$ --- and yet produces
information-rich images whose deflections are sensitive to all of them
simultaneously \citep{refsdal1964,treu2010annurev}. With the
{\it Vera C. Rubin Observatory}\ Legacy Survey of Space and Time (LSST),
the {\it Nancy Grace Roman Space Telescope}, and {\it Euclid}\ now
delivering or imminently delivering wide-field imaging at depth, the
number of known galaxy-scale strong lenses is set to grow from
$\mathcal{O}(10^{3})$ today to $\mathcal{O}(10^{5})$ within a decade.
The bottleneck for cosmological and population-level analyses is no
longer discovery but {\it modeling at scale}.

\gigalens\ \citep{gu2022gigalens} was designed to break that
bottleneck. It expresses the forward lens-model likelihood as a
differentiable JAX/TensorFlow program, exploits hardware parallelism
through \texttt{pmap}/\texttt{shard\_map}, and replaces traditional MCMC
with a hybrid pipeline: stochastic variational inference (SVI) for fast,
initial posterior exploration followed by short Hamiltonian Monte Carlo
chains for refinement. On modern GPU hardware \gigalens\ has been
demonstrated to reduce per-system modeling times from days to minutes,
which is the throughput required to model the LSST/Roman/Euclid lens
catalog within a single graduate-student-year of effort.

The DESI Strong Lens Foundry program was conceived as a follow-up
imaging and modeling campaign for promising lens candidates discovered
in the DESI Legacy Imaging Surveys
\citep{huang2020decals,huang2021desi}. Foundry I
\citep{huang2025foundry1}, the focus of the present paper, is the
first in the series: it presents {\it HST}\ WFC3/F140W observations of
ten high-priority systems (GO--15867; PI Huang) and uses one system,
\foundry, as the methodological demonstration. The Foundry I
\gigalens\ recipe --- an EPL$+$external shear deflector, a S\'{e}rsic
plus shapelet source, and an HMC posterior --- has since been propagated
to other GIGA-Lens-lineage modeling papers, including the
DESI--253.2534$+$26.8843 Einstein-cross analysis of
\citet{cikota2023einsteincross} and the multi-source Carousel-Lens
work of \citet{sheu2024carousel}, as well as forthcoming
Foundry II--V volumes.

Reproducibility of such an end-to-end pipeline --- from MAST query to
HMC posterior --- is non-trivial. The published Foundry I posterior is
a Markov chain in a 74-dimensional space, conditioned on a TinyTim PSF,
a custom central pixel mask, and a non-trivial set of priors
(\citealt{huang2025foundry1} Table~2). A modest discrepancy in any
single element of the recipe can shift the posterior by several formal
$\sigma$ along directions the chain itself was constructed to explore.
For the broader Foundry-series program, and indeed for any future
LSST/Roman/Euclid effort that adopts \gigalens-style differentiable
modeling, an open and {\it independently}\ reproduced version of the
pipeline is a useful complement to the original results: it
exposes which assumptions are load-bearing, which choices are
load-light, and how the published posterior responds to perturbations
that the authors themselves could not run.

In this paper we report such a reproduction. Section~\ref{sec:data}
describes how we re-pull the {\it HST}\ data from MAST. Section
\ref{sec:psf} documents an empirical-PSF pipeline that we find
necessary for the reproduction to converge to the paper's posterior
mode. Section~\ref{sec:model} sets up the lens model and priors.
Section~\ref{sec:inference} describes our inference pipeline, the
HMC compile-time obstacle that forced us to fall back on the SVI
Gaussian variational posterior, and a subsequent v11 NUTS attempt that
resolves the JIT-compile blocker but does not yield a mixed chain.
Section~\ref{sec:results} presents the quantitative comparison to
\citet{huang2025foundry1} Table~3 and addresses the multimodality we
observed in the multi-start MAP search and (separately) in the v11
NUTS run. Section~\ref{sec:discussion} interprets the residual systematic
and itemizes upstream software issues. Section~\ref{sec:software}
documents the software stack. Appendix~\ref{app:ablation} traces the
internal ablation v1\,$\to$\,v11 that led to the working recipe;
Appendix~\ref{app:priors} reproduces the prior table.

%%%%%%%%%%%%%%%%%%%%%%%%%%%%%%%%%%%%%%%%%%%%%%%%%%%%%%%%%%%%%%%%%%%%%%%%%%%%%%%
\section{Data}
\label{sec:data}
%%%%%%%%%%%%%%%%%%%%%%%%%%%%%%%%%%%%%%%%%%%%%%%%%%%%%%%%%%%%%%%%%%%%%%%%%%%%%%%

The Foundry I system \foundry\ was observed with {\it HST}/WFC3-IR in the
F140W filter under program GO--15867 (PI Huang). The combined drizzled
exposure used in this work has a total integration of $1197.7\,$s in
three exposures of $399.23\,$s, was processed through the standard
WFC3 calibration pipeline, and is available from MAST under DOI
10.17909/hx0v-9260. We retrieved the high-level pipeline-combined
drizzle product \texttt{hst\_15867\_65\_wfc3\_ir\_f140w\_ie5065\_drz.fits}
via the \texttt{astroquery.mast.Observations} interface.

The full drizzled science extension is $1523\times1528$ pixels at the
WFC3-IR native pixel scale of $0\farcs13\,\mathrm{pix}^{-1}$. We cut a
$128\times128$ box centered on
$(\alpha,\delta)=(165\fdg4754,-6\fdg0423)$, the source position
given by \citet{huang2025foundry1}. We sky-subtract by the median of
the cutout and adopt a per-pixel noise floor
$\sigma_{\mathrm{bg}}\approx0.029\,\mathrm{e^{-}\,s^{-1}}$ from the
inverse-variance \texttt{WHT} extension. Figure~\ref{fig:cutout} shows
the sky-subtracted science image with an asinh stretch alongside the
\texttt{WHT} extension and a logarithmic display.

\begin{figure*}[!t]
\centering
\includegraphics[width=0.95\textwidth]{figures/cutout_preview.png}
\caption{The $128\times128$ WFC3/F140W cutout of \foundry\ at
$0\farcs13\,\mathrm{pix}^{-1}$. Left to right: sky-subtracted science,
asinh-stretched science, and \texttt{WHT}-derived noise map. The arc
system is well-resolved at native pixel scale; the central deflector
and a fainter companion (modeled as a separate S\'{e}rsic pair; see
\S\ref{sec:model}) are also visible.}
\label{fig:cutout}
\end{figure*}

%%%%%%%%%%%%%%%%%%%%%%%%%%%%%%%%%%%%%%%%%%%%%%%%%%%%%%%%%%%%%%%%%%%%%%%%%%%%%%%
\section{Empirical PSF Construction}
\label{sec:psf}
%%%%%%%%%%%%%%%%%%%%%%%%%%%%%%%%%%%%%%%%%%%%%%%%%%%%%%%%%%%%%%%%%%%%%%%%%%%%%%%

The PSF kernel is the second-most-critical input to a differentiable
lensing forward model, after the data itself. The \gigalens\
distribution ships with a TinyTim WFC3/F140W kernel of shape
$13\times13$ at $0\farcs065\,\mathrm{pix}^{-1}$, i.e. with $\pm
0\farcs42$ support. We found this kernel to be too compact for the
present analysis (\S\ref{sec:results}), and built an empirical PSF
from in-field stars in the same exposure, following the approach
described by \citet{huang2025foundry1}.

The pipeline (script \texttt{17\_build\_empirical\_psf.py}) is:
\begin{enumerate}
\item Open the full $1523\times1528$ drizzled science extension.
Compute sigma-clipped background statistics with
\texttt{astropy.stats.sigma\_clipped\_stats} ($\sigma=3$).
\item Detect bright sources with \texttt{photutils.detection.DAOStarFinder}
at FWHM $=3\,$pix, threshold $50\sigma_{\mathrm{bg}}$, with sharpness
in $[0.2,1.0]$ and roundness in $[-0.5,0.5]$.
\item Reject candidates within $200\,$pix of the lens centroid and
within $25\,$pix of any image edge; sort the remainder by peak flux
and retain the top 30.
\item Pass these to \texttt{photutils.psf.EPSFBuilder}
\citep{bradley2024photutils} with $2\times$ oversampling
($0\farcs065\,\mathrm{pix}^{-1}$ effective), the quartic smoothing
kernel, eight outer iterations, and 20 recentering iterations per
build.
\item Crop the resulting effective PSF to $27\times27$ supersampled
pixels (a tight box of $\pm0\farcs85$ on a side) and normalize to unit
sum.
\end{enumerate}

Figure~\ref{fig:psf} shows our empirical PSF beside the gigalens-default
TinyTim kernel and the median-stacked top-9 raw star postage stamps.
The empirical kernel is broader than TinyTim by roughly a factor of
two in the core, and carries finite power out to the $\pm0\farcs85$
radius --- a region in which TinyTim is identically zero. We will
return to the implications of this width difference in
\S\ref{sec:results}.

\begin{figure*}[!t]
\centering
\includegraphics[width=0.95\textwidth]{figures/psf_comparison.png}
\caption{Empirical PSF construction. Left to right: median of the
top-9 raw F140W star postage stamps; the $27\times27$ supersampled
($0\farcs065\,\mathrm{pix}^{-1}$) empirical PSF built by
\texttt{EPSFBuilder} from 30 isolated stars in the same drizzle;
the TinyTim F140W kernel shipped with \gigalens\ ($13\times13$ at
the same pixel scale); the peak-normalized difference (empirical $-$
TinyTim) with diverging colormap. The TinyTim kernel underestimates
the F140W PSF in the wings beyond $\pm0\farcs42$.}
\label{fig:psf}
\end{figure*}

%%%%%%%%%%%%%%%%%%%%%%%%%%%%%%%%%%%%%%%%%%%%%%%%%%%%%%%%%%%%%%%%%%%%%%%%%%%%%%%
\section{Lens Model}
\label{sec:model}
%%%%%%%%%%%%%%%%%%%%%%%%%%%%%%%%%%%%%%%%%%%%%%%%%%%%%%%%%%%%%%%%%%%%%%%%%%%%%%%

We adopt the same parametric lens model as \citet{huang2025foundry1}:

\subsection{Mass}
A single elliptical power-law convergence
(\texttt{EPL}, \citealt{gu2022gigalens,birrer2018lenstronomy}) with
Einstein radius \thetaE, deprojected 3D mass-density slope \gammaEPL,
ellipticity $(\eone,\etwo)$, and centroid
$(c_{x}^{\mathrm{mass}},c_{y}^{\mathrm{mass}})$,
plus an external shear with two Cartesian components
$(\gextone,\gexttwo)$. This is 8 free mass parameters.

\subsection{Lens light}
Four elliptical S\'{e}rsic profiles with
\texttt{use\_lstsq=False} (i.e.\ the per-component intensity
$I_{e}$ is sampled as a free parameter rather than recovered by an
unconstrained least-squares solver). Two of the four model the central
deflector itself (one compact and one extended); the remaining two
model a fainter companion galaxy detected at $(-2\farcs34,-2\farcs86)$
relative to the main-lens centroid in the v2 residual map
(Figure~\ref{fig:nearby}, see \S\ref{sec:results}). Allowing
$I_{e}$ to be negative through an unconstrained least-squares pass
turned out to be a critical pitfall: an unconstrained fit can drive
some $I_{e}$ negative and absorb $\chi^{2}$ at non-physical
configurations, lowering Einstein-radius accuracy from $0.5\%$ to
$43\%$ (Appendix~\ref{app:ablation}, v4 versus v5).

\subsection{Source light}
A single elliptical S\'{e}rsic plus a Cartesian shapelet basis of
maximum order $n_{\max}=6$ \citep{gu2022gigalens}, giving 28 amplitude
coefficients in addition to the basis-scale parameter $\beta$ and the
shapelet centroid. With \texttt{use\_lstsq=False} we sample the
amplitudes explicitly under a per-coefficient
$\mathcal{N}(0,\,5/\sqrt{i+1})$ prior, with the order-decaying scale
adopted from the canonical \gigalens\ shapelets demonstration. Free-sign
sampling here is appropriate (shapelet basis coefficients can be of
either sign by construction) and complementary to the positive-sign
constraint on lens-light $I_{e}$.

\subsection{Central mask}
\citet{huang2025foundry1} adopt a 2.5-pixel circular mask at the
drizzled $0\farcs065\,\mathrm{pix}^{-1}$ scale to exclude the
central pixels where the lens-light $\rightarrow$ source-light
decomposition is degenerate. Our model lives on the WFC3-IR native
$0\farcs13\,\mathrm{pix}^{-1}$ grid, so we mask the 1.5-pixel circle
($\approx0\farcs20$ radius). The mask is implemented as an inflation
of the per-pixel error map \texttt{err\_map} to $10^{10}$ inside the
mask --- this is the smallest change to \gigalens\,2.0 that yields a
masked likelihood, because the upstream
\texttt{gigalens.jax.model.ForwardProbModel} class does not yet
expose a mask argument. The full custom \texttt{masked\_log\_prob}
function is monkey-patched into the prob model as a JIT-compiled
method.

\subsection{Free-parameter count}
Mass: 8. Lens light (S\'{e}rsic): $4\times7=28$.
Source S\'{e}rsic: 7. Source shapelets: $1+2+28=31$. Total: 74 free
non-linear parameters. The full prior specification is given in
Table~\ref{tab:priors} (Appendix~\ref{app:priors}).

\begin{figure}[!t]
\centering
\includegraphics[width=0.48\textwidth]{figures/nearby_galaxy_detection.png}
\caption{Detection of the nearby galaxy on the v2 residual map.
A clear point-like positive residual at
$(\Delta\alpha,\Delta\delta)\approx(-2\farcs34,-2\farcs86)$ from
the main-lens centroid is the signature of an unmodeled companion;
including its light as two additional S\'{e}rsic components
(v3$\to$v4 onward) removes this feature.}
\label{fig:nearby}
\end{figure}

%%%%%%%%%%%%%%%%%%%%%%%%%%%%%%%%%%%%%%%%%%%%%%%%%%%%%%%%%%%%%%%%%%%%%%%%%%%%%%%
\section{Inference Pipeline}
\label{sec:inference}
%%%%%%%%%%%%%%%%%%%%%%%%%%%%%%%%%%%%%%%%%%%%%%%%%%%%%%%%%%%%%%%%%%%%%%%%%%%%%%%

\subsection{Hardware}
All fits ran on \texttt{/raid/benson}, an AlmaLinux 9.7 aarch64 host
with two NVIDIA L4 (23\,GB) and eight NVIDIA A16 (15\,GB) GPUs visible
to CUDA. \gigalens's \texttt{shard\_map}-parallelised
\texttt{ModellingSequence.MAP} and \texttt{.SVI} routines build their
compiled executables for a single GPU type and then attempt to dispatch
across all visible devices, which fails on this host with a
\texttt{cannot run on CUDA:2 of type 'NVIDIA A16'} error. We work
around this by UUID-pinning the two L4s via
\texttt{CUDA\_VISIBLE\_DEVICES=GPU-$\langle$uuid-of-L4-0$\rangle$,GPU-$\langle$uuid-of-L4-1$\rangle$}.
Once pinned, JIT compilation succeeds and the two L4s deliver
$\sim$1.7 steps/s on the 200-chain, 400-step v9 MAP problem
(\S\ref{sec:results}), or $\sim$3 minutes of wall time after JIT.

\subsection{MAP search (Stage A)}
We run a 200-chain Adabelief \citep[][lr$=10^{-2}$, $b_{1}=0.95$,
$b_{2}=0.99$]{birrer2018lenstronomy} multi-start MAP search over
400 steps. Each chain is seeded by an independent prior sample.
\texttt{ModellingSequence.MAP(\dots,\,output\_type='best')} returns the
single chain with the highest log posterior at the final step. Crucially
we also retain the full
\texttt{(n\_chains,\,n\_steps,\,n\_params)} array so that we can
characterize the multimodality of the posterior surface; this is the
v7/v9 multi-start analysis of \S\ref{sec:results}.

\subsection{SVI (Stage B)}
Hamiltonian Monte Carlo was attempted but did not complete in any
reasonable wall time. Specifically, the JIT compile of
\texttt{pmap(scan(mcmc\_sample\_chain(PreconditionedHMC} $+$
\texttt{GradientBasedTrajectoryLengthAdaptation)))} did not finish
within several hours of compile time on the present 74-parameter model
under JAX 0.6.2 \citep{bradbury2018jax}. We did not pursue this
further; for the present paper we adopt a Gaussian stochastic variational
posterior as a surrogate. This is documented honestly throughout
(\S\ref{sec:discussion}): we are not running HMC, we are running SVI,
and our $1\sigma$ widths in Table~\ref{tab:mass} are
{\it variational}-posterior widths (and therefore underestimates of
the marginal posterior width if the true posterior is non-Gaussian
in the parameters of interest).

The SVI is initialised at the best paper-mode chain found by the
preceding 200-chain MAP search (\texttt{map\_v9\_paper\_mode.npz}),
with $n_{\mathrm{VI}}=200$ particles, an Adabelief learning rate of
$10^{-4}$, init scales of $10^{-4}$, and 500 steps. The tight
initialisation and the small learning rate are intentional: they pin
the variational posterior to the paper-mode local minimum rather than
allowing it to drift to one of the higher-$\chi^{2}$ but lower-log-$p$
modes that we identified in v7 (\S\ref{sec:results}).

After 500 SVI steps the negative ELBO has fallen monotonically from
$5.69\times10^{4}$ to $4.58\times10^{4}$; the chain shows no sign of
mode-hopping, and the final variational mean is essentially indistinguishable
from the initial MAP. We then draw $10^{4}$ samples from $q(z)$ and
push them through the bijector to recover the physical-parameter
posterior used in Table~\ref{tab:mass}.

\subsection{Wall time}
Total wall time for the full v9 multi-start MAP plus v10 SVI is
$\approx7\,$min ($3.9\,$min for MAP, $4.9\,$min for SVI), excluding
the one-time JIT compile (which dominates the first run; for the SVI
specifically, cuDNN convolution autotuning takes $\sim 90\,$s due to
the $f32[100,1,256,256]$ convolution shape required by the $74$-parameter
model on 2 GPUs).

\subsection{NUTS attempt (v11 series)}
\label{sec:nuts}

Having identified the HMC JIT compile-time blowup as the most
load-bearing of our upstream issues (\S\ref{sec:discussion}), we
attempted a sequence of six NUTS runs (v11, v11b, v11c, v11d-scan,
v11e, v11f) targeting that issue specifically. We treat the resulting
chain as a {\it diagnostic}, not a replacement headline posterior;
the v10 SVI fit of \S\ref{sec:headline} remains the canonical result.

\paragraph{NUTS does compile.} The replacement we adopt is
\texttt{tfp.experimental.mcmc.PreconditionedNoUTurnSampler} (TFP\,0.25.0,
JAX backend) with \texttt{max\_tree\_depth=6} pinned as a static argument
and {\it without}\ a \texttt{pmap} wrapper. With these two choices the
sampler JIT-compiles in $\sim$5\,min on the same 2$\times$L4 host that
failed to compile HMC after several hours, and runs 700 chain steps
(200 burn-in $+$ 500 kept) in $\sim$2--3\,min of wall time. This is
the main positive finding of the v11 series: the JIT-tractability of
sample-chain at $74$ parameters under JAX\,0.6.2 is recovered by NUTS
$+$ static tree depth, even though the same combinatorics defeated
\texttt{PreconditionedHMC}$+$\texttt{Gra\-dient\-Based\-Tra\-jec\-tory\-Length\-Adap\-ta\-tion}.

\paragraph{The v10 SVI covariance is rank-deficient.} On loading
\texttt{svi\_v10\_posterior\_mass.npz} and computing the SVI posterior
covariance in $z$-space, we find that \texttt{np.linalg.eigvalsh(qz\_cov)}
reports a minimum eigenvalue of $-4.4\times10^{-13}$ --- numerically zero
but slightly negative --- and a numerical rank of $53/74$
(\texttt{np.linalg.matrix\_rank}, default tolerance), implying $\sim 21$
near-degenerate directions. The CPU NumPy Cholesky
(\texttt{np.linalg.cholesky}) raises \texttt{NotPositiveDefinite}; the
JAX Cholesky (\texttt{jnp.linalg.cholesky}) {\it silently returns NaN}.
The first NUTS attempt (v11) used this raw covariance as the
\texttt{precondition\_state\_param\_partial\_loc} and inherited NaN
momentum samples; TFP's NUTS proceeds without erroring and produces
near-zero proposals. The fix is to jitter the covariance by
$10^{-6}\,\mathbb{I}$ before Cholesky, after which the Cholesky
diagonal lies in $[1.0\times10^{-3},6.2\times10^{-2}]$ and is usable.

\paragraph{Dual-averaging step-size collapse.} Even with a valid
(jittered) preconditioner (v11e), TFP's
\texttt{DualAveragingStepSizeAdaptation} consistently drove the step
size to $\sim$2--$3\times10^{-6}$ by the end of the 200-step burn-in,
at which point every proposed move is effectively the current state
(acceptance probability $1.000$, leapfrog count clamped at the
$2^{6}-1=63$ maximum, and per-step displacement $\|\Delta z\|\sim
10^{-4}$--$10^{-3}$ in unconstrained space). A static step-size scan
over $\{10^{-5},3\times10^{-5},10^{-4},3\times10^{-4},10^{-3},
3\times10^{-3},10^{-2},3\times10^{-2}\}$ with adaptation disabled
(v11d-scan) gave 0\% acceptance for every entry $\ge10^{-5}$ --- the
chain immediately U-turns and rejects --- and reported
\texttt{leapfrogs\_taken=1} at every iteration. A working step size
therefore exists somewhere in the interval $(2\times10^{-6},10^{-5})$,
but neither simple dual-averaging nor a coarse log-spaced scan locates
it; the local Hessian along the SVI mode and the local Hessian along
the basin we actually want to explore appear to differ enough that a
single global mass matrix is not adequate.

\paragraph{A third local mode (v11f).} As a final attempt, we
re-initialized NUTS from the v7 paper-mode MAP chain
(\texttt{map\_v7\_paper\_mode.npz}; $\log p = -21{,}789$ under TinyTim
PSF, but $\log p = -99{,}119$ under the empirical PSF used here)
rather than the v10 SVI mean, kept the $10^{-6}\,\mathbb{I}$ jitter,
and loosened the dual-averaging target acceptance from the default
$0.8$ to $0.6$. The chain moves $\sim 35{,}000$ $\log p$ units during
burn-in into a {\it new}\ local mode with $\log p \approx -63{,}868$,
distinct from both the v10 SVI mode ($\log p \approx -40{,}400$) and
the original paper-mode MAP. The posterior medians at this new mode
are summarized in Table~\ref{tab:nuts_v11f}. With $\thetaE=2.604''$,
v11f reproduces the published Einstein radius to $1.6\%$ --- the
closest of any version of this reproduction --- although the other
five mass parameters move in mixed directions relative to the v10
medians. We emphasize that v11f is a {\it converged-to-a-different-basin}\
point estimate, not a mixed chain; its $1\sigma$ widths
(Table~\ref{tab:nuts_v11f}) are widths of a chain that froze and are
not faithful marginal-posterior widths.

\paragraph{Why we did not pursue \texttt{windowed\_adaptive\_nuts}.}
The natural next step is TFP's \texttt{tfp.experimental.mcmc.
windowed\_adaptive\_nuts}, which runs Stan-style window-adaptive mass
matrix estimation interleaved with step-size adaptation and would in
principle resolve the curvature mismatch between modes. The blocker
is interface: \texttt{windowed\_adaptive\_nuts} expects a
\texttt{tfd.JointDistribution\-Auto\-Batched} or equivalent, while the
\gigalens\ \texttt{ForwardProbModel} exposes
\texttt{log\_prob}/\texttt{sample} but not the joint-distribution
abstraction. Refactoring the prob model to expose a
\texttt{JointDistribution} interface is straightforward in principle
but well-scoped enough to defer to a follow-up paper; the present
work documents the obstacle so that downstream users know what would
be required.

%%%%%%%%%%%%%%%%%%%%%%%%%%%%%%%%%%%%%%%%%%%%%%%%%%%%%%%%%%%%%%%%%%%%%%%%%%%%%%%
\section{Results}
\label{sec:results}
%%%%%%%%%%%%%%%%%%%%%%%%%%%%%%%%%%%%%%%%%%%%%%%%%%%%%%%%%%%%%%%%%%%%%%%%%%%%%%%

\subsection{Headline mass-parameter reproduction}
\label{sec:headline}

Table~\ref{tab:mass} compares our v10 SVI posterior medians and 1-$\sigma$
percentile widths to the published HMC posterior of
\citet{huang2025foundry1} (their Table~3). The medians and 16/84
percentiles are computed by loading our 10\,000-sample posterior
(\texttt{data/svi\_v10\_posterior\_mass.npz}) and applying
\texttt{numpy.median} and \texttt{numpy.percentile}.

\begin{deluxetable*}{lcccccc}
\tablecaption{Headline mass-parameter reproduction. v10 SVI Gaussian
variational posterior (this work) versus the HMC posterior of
\citet[][Table 3]{huang2025foundry1}.\label{tab:mass}}
\tablehead{
\colhead{Param} & \colhead{v10 median} & \colhead{v10 $1\sigma$ width} &
\colhead{Huang+2025a} & \colhead{Huang+2025a $1\sigma$} &
\colhead{$\Delta_{\mathrm{med}}$} &
\colhead{\% off, sign}
}
\startdata
\thetaE\ (arcsec) & $+2.5659$ & $\pm0.0003$ & $+2.6463$ & $\pm0.0017$ & $-0.0804$ & $3.0\%$, $\checkmark$ \\
\gammaEPL\        & $+2.1507$ & $\pm0.0004$ & $+1.372$  & $\pm0.023$  & $+0.7787$ & $56.8\%$, $\checkmark$ (both $>1$) \\
\eone\            & $+0.1064$ & $\pm0.0004$ & $+0.1091$ & $\pm0.0020$ & $-0.0027$ & $2.5\%$, $\checkmark$ \\
\etwo\            & $-0.0533$ & $\pm0.0002$ & $-0.1320$ & $\pm0.0020$ & $+0.0787$ & $59.6\%$, $\checkmark$ \\
\gextone\         & $+0.0399$ & $\pm0.0003$ & $+0.0657$ & $\pm0.0024$ & $-0.0258$ & $39.3\%$, $\checkmark$ \\
\gexttwo\         & $-0.0676$ & $\pm0.0002$ & $-0.0939$ & $\pm0.0022$ & $+0.0263$ & $28.0\%$, $\checkmark$ \\
\enddata
\tablecomments{Sign-quadrant agreement on all six mass parameters.
\thetaE\ and \eone\ are reproduced to $\le3\%$ of the published median.
The slope \gammaEPL\ converges to a higher local minimum than the
published value, as discussed in \S\ref{sec:multimodal} and
\S\ref{sec:discussion}.
The 1-$\sigma$ widths reported in column 3 are widths of the
{\it variational}\ Gaussian posterior; they underestimate the
marginal posterior width if the latter is non-Gaussian.}
\end{deluxetable*}

In polar form the lens ellipticity vectors are
$|\varepsilon|=0.118$ at position angle $-13^{\circ}$ for our v10
posterior versus $|\varepsilon|=0.171$ at $-25^{\circ}$ for the
published HMC posterior --- a $\sim 30\%$ magnitude offset and a
$12^{\circ}$ residual rotation. The external shears are
$|\gext|=0.079$ at $-29^{\circ}$ versus $|\gext|=0.115$ at $-28^{\circ}$:
the position angle of the shear matches to within $1^{\circ}$, while
the magnitude is $\approx 30\%$ low. These coherent offsets are
revisited in \S\ref{sec:single-constraint}.

\subsection{Second-mode point estimate from v11f NUTS}
\label{sec:nuts_v11f_results}

Table~\ref{tab:nuts_v11f} reports the posterior medians of the v11f
NUTS chain (\S\ref{sec:nuts}; \texttt{data/nuts\_v11f\_posterior\_mass.npz}),
the only NUTS variant that converged to a distinct basin rather than
freezing at the SVI mode. This is presented as a {\it second-mode
point estimate}, not as a replacement for the v10 SVI posterior of
Table~\ref{tab:mass}: v11f's chain froze in its target basin
($\log p\approx-63{,}868$) and did not mix, so its quoted $1\sigma$
widths are widths of a non-mixed chain and do not represent marginal
posterior uncertainty. The Einstein radius at the v11f mode,
$\thetaE=2.604''$, is $1.6\%$ off the published value --- the closest
single-version reproduction of \thetaE\ that we have obtained.

\begin{deluxetable}{lccc}
\tablecaption{v11f NUTS second-mode point estimate.\label{tab:nuts_v11f}}
\tablehead{
\colhead{Param} & \colhead{v11f median} & \colhead{Huang+2025a} & \colhead{\% off}
}
\startdata
\thetaE\ (arcsec) & $+2.6042$ & $+2.6463$ & $1.6\%$ \\
\gammaEPL\        & $+2.2565$ & $+1.372$  & $64.5\%$ \\
\eone\            & $+0.1592$ & $+0.1091$ & $45.9\%$ \\
\etwo\            & $-0.1507$ & $-0.1320$ & $14.2\%$ \\
\gextone\         & $+0.0543$ & $+0.0657$ & $17.3\%$ \\
\gexttwo\         & $-0.1186$ & $-0.0939$ & $26.3\%$ \\
\enddata
\tablecomments{v11f is a NUTS chain initialised at the v7 paper-mode
MAP point with a $10^{-6}\,\mathbb{I}$-jittered v10 SVI covariance as
preconditioner and target acceptance $0.6$. The chain converged into
a new local basin $\log p\approx-63{,}868$ (distinct from the v10
SVI mode at $\log p\approx-40{,}400$) and remained there. The
quoted medians are over 500 retained samples after 200 burn-in steps;
$1\sigma$ widths are not reported because the chain did not mix.
v11f reproduces \thetaE\ to $1.6\%$ but is more discrepant than v10
in \gammaEPL\ and \eone. See \S\ref{sec:nuts} for the run details and
\S\ref{sec:multimodal_v11} for the multimodality implication.}
\end{deluxetable}

\subsection{Residual maps}
\label{sec:residuals}

Figure~\ref{fig:residv2} shows the v2-stage residual on which the
nearby galaxy was first detected, and Figure~\ref{fig:residv4} shows
the v4-stage residual --- a cautionary example. The v4 fit
($\chi^{2}=3.45$ on the masked likelihood) appears at first glance to
be a better fit than v5--v6 ($\chi^{2}=7.56$ and $6.59$ respectively),
but the v4 lens-light \texttt{use\_lstsq=True} configuration is
exploiting the freedom to set some component intensities negative
(\S\ref{sec:model}), and its Einstein radius is off by $43\%$
($\thetaE^{\mathrm{v4}}=1.52$). This is the canonical pathology of
unconstrained linear amplitude solvers in a lensing forward model:
$\chi^{2}$ is not the right diagnostic for parameter accuracy if the
model is allowed to absorb residuals through unphysical degrees of
freedom.

We do not have a v10 best-fit residual map available; constructing
one requires re-running the simulator forward pass with the
$10^{4}$ posterior samples, which is out of scope for the present
paper. A v10 residual will be included in supplementary materials.

\begin{figure}[!t]
\centering
\includegraphics[width=0.48\textwidth]{figures/map_v2_residual.png}
\caption{v2 MAP residual (data $-$ model, in units of
$\sigma_{\mathrm{bg}}$). The bright residual blob $\sim3''$ southwest
of the main lens is the signature of the nearby galaxy modeled in
v3 onward (Figure~\ref{fig:nearby}).}
\label{fig:residv2}
\end{figure}

\begin{figure}[!t]
\centering
\includegraphics[width=0.48\textwidth]{figures/map_v4_residual.png}
\caption{v4 MAP residual. Despite an apparently excellent
$\chi^{2}=3.45$, this fit places the Einstein radius at
$1\farcs52$, $43\%$ low of the published value, because the
unconstrained least-squares lens-light intensity solver is allowed
to drive some S\'{e}rsic amplitudes negative. Lesson: positive
amplitudes are a non-negotiable physical constraint that must be
enforced in the prior rather than the loss.}
\label{fig:residv4}
\end{figure}

\subsection{Multimodality of the posterior surface}
\label{sec:multimodal}

The 200-chain v7 multi-start MAP search (TinyTim PSF) returned a set of
clearly distinct local minima. Of the 200 random-prior starts:
\begin{itemize}
\item $17/200 = 8.5\%$ converged to the paper's quadrant
($\eone>0$, $\etwo<0$, $\gextone>0$, $\gexttwo<0$); the best
such chain (idx 140) reached $\chi^{2}=7.72$, $\log p=-2.18\times10^{4}$.
\item $106/200 = 53\%$ had $\etwo<0$ (any sign on the others); the
unconstrained global best ($\chi^{2}=6.67$, $\log p=-1.33\times10^{4}$;
chain idx unspecified at this stage) had $\etwo>0$ and was therefore
{\it not}\ in the paper's quadrant.
\item The top three chains by $\log p$ differed not just in
parameter values but in the {\it sign quadrant}\ of $(\etwo,\gexttwo)$:
the same physical system admits at least two qualitatively distinct
mass-model orientations consistent with the data plus our priors.
\end{itemize}
The paper's chain is therefore not the global best under our
priors and PSF, even though our priors closely follow those of
\citet[][Table 2]{huang2025foundry1}: see Appendix~\ref{app:priors}.

\subsection{Empirical-PSF flip}
\label{sec:psfflip}

Replacing TinyTim with the empirical PSF of \S\ref{sec:psf} produced
a striking change in the topology of the multi-start search. With the
empirical kernel, in the v9 200-chain search:
\begin{itemize}
\item paper-mode chain fraction increased $17/200 \to 45/200$, i.e.\
$8.5\%\to 22.5\%$;
\item the fraction of chains with $\gammaEPL<1.7$ (the paper sits at
$\gammaEPL=1.37$) increased from $\sim 5\%$ to $35\%$
($70/200$ in v9);
\item the {\it global best}\ chain (idx 163, $\chi^{2}=11.93$,
$\log p=-5.64\times10^{4}$) is itself a paper-mode chain.
\end{itemize}
The $\chi^{2}$ floor under the empirical PSF rose from 6.67 to 11.93,
which is what we should expect: TinyTim is over-smoothed relative to
the true F140W PSF (Figure~\ref{fig:psf}), so the model had spurious
freedom to ``earn'' $\chi^{2}$ improvement by absorbing residuals into
the source shapelets. Forcing the model to use a wider, more honest
PSF raises the achievable $\chi^{2}$ floor and re-centers the
posterior on a physically meaningful mode.

The v10 SVI posterior reported in Table~\ref{tab:mass} is the
$10^{4}$-sample Gaussian variational posterior pinned to this v9
global-best paper-mode chain.

%%%%%%%%%%%%%%%%%%%%%%%%%%%%%%%%%%%%%%%%%%%%%%%%%%%%%%%%%%%%%%%%%%%%%%%%%%%%%%%
\section{Discussion}
\label{sec:discussion}
%%%%%%%%%%%%%%%%%%%%%%%%%%%%%%%%%%%%%%%%%%%%%%%%%%%%%%%%%%%%%%%%%%%%%%%%%%%%%%%

\subsection{Coherent residual signature suggests a single missing
constraint}
\label{sec:single-constraint}

The residuals between our v10 posterior and \citet{huang2025foundry1}
Table 3 are not random across the six parameters --- they are coherent
along directions that suggest one underlying physical effect. Specifically:
\begin{itemize}
\item The lens ellipticity vector rotates by $12^{\circ}$
(\eone\ accurate to $2.5\%$, but \etwo\ is $\sim 60\%$ low in magnitude
in a sign-consistent way).
\item The external shear position angle matches to within $1^{\circ}$,
but the magnitude is $\sim 30\%$ low.
\item The deprojected slope \gammaEPL\ runs to a value of 2.15 instead
of the published 1.37, a 57\% increase in slope; in our priors
$\gammaEPL$ has a lower bound at 1.0, and the chain finds a stable
local minimum well above it. \citet{huang2025foundry1} themselves
note (their \S 4) that the published posterior was ``pushing
against the prior lower bound'' --- their reported median sits at
$\gammaEPL=1.37$ with $1\sigma=0.023$, but the prior range starts at
$1.0$, leaving little room.
\end{itemize}
Each of these three effects --- ellipticity-vector rotation, smaller
shear magnitude, higher slope --- is consistent with a model in which
the lens-light and mass-model position angles are not tied, so the
mass model is free to adopt a slightly different orientation than
the lens light and partially absorb the missing constraint into a
mix of $|\etwo|$, $|\gext|$, and \gammaEPL. The paper's mass--light PA
tying (or an equivalent multipole-correction term) would in principle
pin the lens ellipticity to the lens-light position angle and remove
this rotational freedom, reducing the model to a sub-region of our
parameter space where the slope can drift to the lower-\gammaEPL\
mode found by the published HMC.

Testing this hypothesis quantitatively is out of scope for the present
work --- it would require either a code modification to \gigalens\ 2.0
(to expose a mass--light PA tying argument) or a fully custom log-
likelihood that imposes the tying as a soft penalty. Both are natural
next steps for v11 of this reproduction.

\subsection{The PSF is a load-bearing assumption}

Sections~\ref{sec:psf} and \ref{sec:psfflip} together demonstrate that
the choice of PSF is not a cosmetic detail of the pipeline. Replacing
TinyTim with the empirical kernel from in-field stars flipped the
global-best chain into the paper's mode. The mechanism is plausible:
a too-narrow PSF gives the source-shapelet basis more freedom to
absorb central residuals, lowering the model's marginal preference
for a steep ($\gammaEPL\gtrsim 2$) deflector, but only in modes where
the ellipticity vector accommodates the spurious flux. Any future
{\it HST}\ \gigalens\ reproduction --- including the still-pending
Foundry II--V volumes --- should build an empirical PSF from field
stars in the same exposure before fitting.

\subsection{Multimodality is a real feature of the posterior surface}
\label{sec:multimodal_v11}

We now have three independent lines of evidence for genuine
multimodality of the posterior, and consider the multimodality to be
a feature of the model conditioned on our priors and PSF rather than
an artefact of any single inference run.

\paragraph{Evidence (a): v7 multi-start MAP.} The 200-chain v7
multi-start search revealed at least three distinct local minima
differing in the sign quadrant of $(\etwo,\gexttwo)$ at fixed positive
\eone\ and positive \gextone\ (\S\ref{sec:multimodal}); only $17/200
=8.5\%$ of chains land in the paper's quadrant under TinyTim.

\paragraph{Evidence (b): v6 vs v8 mode swap.} The v6 single-chain MAP
(TinyTim, $\chi^{2}=6.59$) sits in a sign quadrant
$(\eone,\etwo,\gextone,\gexttwo)=(+,+,+,+)$ that is different from the
v8 SVI mode $(+,-,+,-)$, even though both converged stably under the
same loss function (Appendix~\ref{app:ablation}). The two are
qualitatively distinct lensing geometries, not refinements of one
another.

\paragraph{Evidence (c): v11f NUTS basin transit.} The v11f NUTS run
(\S\ref{sec:nuts}, \S\ref{sec:nuts_v11f_results}) was initialised at
the v7 paper-mode MAP point and during burn-in moved $\sim 35{,}000$
$\log p$ units into a basin at $\log p\approx -63{,}868$, distinct
from the v10 SVI basin at $\log p\approx -40{,}400$. The medians
at v11f's terminal basin reproduce $\thetaE$ to $1.6\%$ but have a
different signature in $(\eone,\etwo,\gext)$ than v10 does --- in
particular $\eone=+0.159$ (vs.\ v10's $+0.106$) and $\etwo=-0.151$
(vs.\ v10's $-0.053$). This is a third distinct local mode, found
by a gradient-based sampler rather than by the multi-start scan, and
confirms that the model surface admits at least three usable basins
under our priors and PSF.

The published HMC posterior is centered in only one of these basins.
Two interpretations are possible: (i) the published HMC chain explored
other modes but with very low weight, and the reported median is
correctly centered; (ii) the published HMC chain was started in the
paper-mode basin and remained there for the duration of the run.
We cannot distinguish these from the published Foundry I paper alone,
and neither is unreasonable. The implication, however, is the same:
when running \gigalens\ on a new system without an external prior on
ellipticity sign, a multi-start MAP scan {\it followed by}\
mode-aware sampling (or, in lieu of that, transparent reporting of
which local basin the chain occupies) is the right first step ---
gradient-based local search returns a posterior mode, not the
posterior.

\subsection{Upstream issues encountered}
\label{sec:upstream}

We encountered six issues in the \gigalens\,2.0 / JAX\,0.6.2 /
TFP\,0.25.0 stack that we believe are appropriate to file as GitHub
issues:
\begin{enumerate}
\item \textbf{Implicit \texttt{shard\_map} import.} \gigalens\,2.0
references \texttt{jax.experimental.shard\_map.shard\_map} by attribute
without first ensuring the submodule is imported. Under JAX 0.6.2 this
fails with \texttt{AttributeError: module 'jax.experimental' has no
attribute 'shard\_map'} unless the caller pre-imports
\texttt{jax.experimental.shard\_map}. We work around this by adding
\texttt{import jax.experimental.shard\_map} to every fitting script.
\item \textbf{Negative lens-light amplitudes via unconstrained
\texttt{lstsq}.} The \gigalens\ simulator path that resolves S\'{e}rsic
intensities via \texttt{jnp.linalg.lstsq} can return negative
amplitudes. These are unphysical, and the resulting MAP can absorb
$\chi^{2}$ at non-physical configurations (Appendix~\ref{app:ablation},
v4 versus v5). A NNLS-based variant or a clipped output would resolve
this. We disable this code path entirely by setting
\texttt{use\_lstsq=False} on every S\'{e}rsic component.
\item \textbf{HMC compile time at $74$ parameters.} The combination
\texttt{pmap}$\circ$\texttt{scan}$\circ$\texttt{mcmc\_sample\_chain}
$\circ$\texttt{PreconditionedHMC}$+$
\texttt{GradientBasedTrajectoryLengthAdaptation} did not finish
JIT-compiling in any wall time we were willing to wait on this
$74$-parameter model. Switching to
\texttt{tfp.experimental.mcmc.PreconditionedNoUTurnSampler} with a
static \texttt{max\_tree\_depth} and no \texttt{pmap} wrapper
(\S\ref{sec:nuts}) compiles in $\sim$5\,min and runs in $\sim$2--3\,min,
which resolves the JIT-compile blocker. The published HMC posterior
must, however, have been run on a different combination of versions
(JAX, gigalens branch); a standalone reproducer of the original HMC
compile-time blowup would still be useful for the JAX/TFP issue tracker.
\item \textbf{Mixed-GPU-type \texttt{shard\_map} crash.} On heterogeneous
hosts (e.g.\ our 8$\times$A16 $+$ 2$\times$L4), the executable
compiled for the first sharded device cannot be run on the others, and
the shard map crashes mid-run with
\texttt{executable is built for device CUDA:0 of type 'NVIDIA L4'; cannot
run on CUDA:2 of type 'NVIDIA A16'}. UUID-pinning the visible devices
is a workaround; a cleaner fix would be for the shard map machinery
to recompile per-device or to refuse to shard across heterogeneous
device types with an actionable error.
\item \textbf{\texttt{jnp.linalg.cholesky} silently returns NaN on
a rank-deficient SVI covariance.} The v10 SVI posterior covariance
has $\min\lambda(qz\_cov)=-4.4\times10^{-13}$ and numerical rank
$53/74$ (\S\ref{sec:nuts}). On this matrix the CPU NumPy Cholesky
raises \texttt{NotPositiveDefinite}, but the JAX Cholesky silently
returns NaN. NUTS then receives NaN momentum and proceeds without
erroring, producing the appearance of a stuck chain rather than a
diagnosable failure. At minimum, \texttt{jnp.linalg.cholesky} should
emit a NaN-trace warning under \texttt{jax\_debug\_nans} the way other
NaN-producing JAX primitives do. The downstream \gigalens\ pipeline
would benefit from a defensive \texttt{cov + eps*I} jitter in any code
path that feeds an SVI covariance into a TFP preconditioner.
\item \textbf{\texttt{DualAveragingStepSizeAdaptation} pathology on
multimodal-Hessian models.} On this 74-parameter model the dual-averaging
adapter consistently drove the step size to $\sim$2--$3\times10^{-6}$
within 200 burn-in steps, at which acceptance is $1.000$ but per-step
$\|\Delta z\|$ is $10^{-4}$--$10^{-3}$ and the chain effectively does
not move (\S\ref{sec:nuts}). A static log-spaced scan from $10^{-5}$
to $3\times10^{-2}$ with adaptation disabled gives $0\%$ acceptance
everywhere, so the working step size lives in the narrow window
$(2\times10^{-6},10^{-5})$ --- below the resolution of a coarse scan
and below where dual-averaging settles. We interpret this as the
adapter optimizing a single global step-size for a target acceptance
under a globally Gaussian-like assumption that the model surface
violates: the local Hessian at the v10 SVI mode and the local Hessian
at the v11f basin differ enough that no single mass matrix or step
size is simultaneously appropriate. A windowed mass-matrix adapter
(e.g.\ \texttt{windowed\_adaptive\_nuts}) or a per-basin re-tuning
would in principle help; both require a \texttt{JointDistribution}
refactor of \gigalens\ which we defer.
\end{enumerate}

\subsection{What full reproduction would require}

To reproduce the published HMC posterior bit-for-bit one would need
(i) the paper's exact \gigalens\ git revision (the multi-node branch
of \texttt{giga-lens/gigalens} has moved since the paper was written);
(ii) mass--light position-angle tying, which is not exposed in the
\gigalens\,2.0 public API as of our run date and would require a
modification to the simulator internals; (iii) optionally a
higher-order multipole correction term in the EPL profile;
(iv) a working HMC at $74$ parameters in JAX 0.6.2, which we did not
manage to obtain (NUTS with static \texttt{max\_tree\_depth} resolves
the JIT-compile blocker, \S\ref{sec:nuts}, but step-size adaptation
on this multimodal-Hessian surface remains an open problem); and
(v) a \texttt{JointDistribution} refactor of the \gigalens\ prob model
to enable \texttt{windowed\_adaptive\_nuts} or similar window-adaptive
mass-matrix samplers. These are the natural next-steps for any reader
who wishes to push this reproduction further.

%%%%%%%%%%%%%%%%%%%%%%%%%%%%%%%%%%%%%%%%%%%%%%%%%%%%%%%%%%%%%%%%%%%%%%%%%%%%%%%
\section{Software and Data Availability}
\label{sec:software}
%%%%%%%%%%%%%%%%%%%%%%%%%%%%%%%%%%%%%%%%%%%%%%%%%%%%%%%%%%%%%%%%%%%%%%%%%%%%%%%

The \emph{HST}\ data used in this paper is public via MAST
DOI 10.17909/hx0v-9260. The \gigalens\ source is at
\url{https://github.com/giga-lens/gigalens} (the \texttt{multi-node}
branch at commit-time was the working baseline for this paper).
The reproduction scripts (\texttt{01\_download\_hst.py} through
\texttt{19\_svi\_v10\_paper\_mode\_empirical.py} for the v1$\to$v10
pipeline, and \texttt{20\_fit\_nuts\_v11.py},
\texttt{21\_fit\_nuts\_v11b.py}, \texttt{22\_fit\_nuts\_v11c.py},
\texttt{23\_nuts\_stepsize\_scan.py}, \texttt{24\_fit\_nuts\_v11e.py},
\texttt{25\_fit\_nuts\_v11f.py} for the v11 NUTS series), together
with the analysis notebooks and the per-version posterior artifacts
(\texttt{data/svi\_v10\_posterior\_mass.npz},
\texttt{data/nuts\_v11\{,b,c,e,f\}\_samples.npz}, and the matching
\texttt{nuts\_v11*\_posterior\_mass.npz} files), are available on
request and are planned for public release as a
\texttt{foundry-i-reproduction} repository.

Table~\ref{tab:software} lists the exact pinned versions of the
software stack on the run host.

\begin{deluxetable}{ll}
\tablecaption{Software stack versions.\label{tab:software}}
\tablehead{\colhead{Package} & \colhead{Version}}
\startdata
Python                          & 3.13 \\
JAX                             & 0.6.2 (with jax[cuda12]) \\
TensorFlow                      & 2.20.0 \\
TensorFlow Probability          & 0.25.0 \\
\gigalens                       & 1.2.1 (editable; multi-node branch) \\
\lenstronomy                    & 1.14.0 \\
\photutils                      & 3.0.0 \\
astroquery                      & 0.4.11 \\
\textsc{optax}                  & 0.2.8 \\
\textsc{NumPy}                  & 2.4.6 \\
\textsc{Astropy}                & current LTS \\
CUDA                            & 12 (host runtime 13.0 in PATH) \\
Operating system                & AlmaLinux 9.7 (aarch64) \\
GPUs used                       & 2$\times$NVIDIA L4 (UUID-pinned) \\
\enddata
\tablecomments{HMC at the 74-parameter model size was not obtained
under this stack; we use the v10 Gaussian SVI as the headline
variational-posterior surrogate (Table~\ref{tab:mass}). NUTS at the
same model size {\it does}\ compile under this stack with a static
\texttt{max\_tree\_depth} (\S\ref{sec:nuts}); v11f provides a
second-mode point estimate (Table~\ref{tab:nuts_v11f}). See
\S\ref{sec:upstream} for the full upstream-issues list.}
\end{deluxetable}

%%%%%%%%%%%%%%%%%%%%%%%%%%%%%%%%%%%%%%%%%%%%%%%%%%%%%%%%%%%%%%%%%%%%%%%%%%%%%%%
\section{Conclusions}
\label{sec:conclusions}
%%%%%%%%%%%%%%%%%%%%%%%%%%%%%%%%%%%%%%%%%%%%%%%%%%%%%%%%%%%%%%%%%%%%%%%%%%%%%%%

We have presented an end-to-end public-data reproduction of the
\gigalens\ modeling demonstration of Foundry I \citep{huang2025foundry1}
on the system \foundry. The principal findings are:
\begin{enumerate}
\item All six mass-model sign quadrants of the published HMC posterior
are reproduced by our SVI Gaussian variational posterior.
\item \thetaE\ matches the published median to $3.0\%$; \eone\ to
$2.5\%$; the external-shear position angle to within $1^{\circ}$.
\item The remaining gap --- a $12^{\circ}$ residual rotation of the
lens ellipticity, $\sim30\%$ smaller $|\etwo|$ and $|\gext|$, and a
$57\%$ steeper \gammaEPL\ --- is coherent rather than scattered, and
is consistent with one missing physical constraint (most likely
mass--light PA tying or a higher-multipole correction) rather than
three independent failures.
\item The Foundry I posterior surface is genuinely multimodal under
the published priors: in 200 random-start MAP fits with TinyTim PSF,
only $8.5\%$ of chains land in the paper's quadrant, and the
unconstrained global best lies in a different ($\etwo>0$) mode.
\item Replacing the gigalens-default TinyTim PSF with an empirical
PSF built from in-field stars flips the global-best chain into the
paper's quadrant ($45/200=22.5\%$, with global best in mode) at the
cost of a higher but more honest $\chi^{2}$ floor.
\item We resolve the upstream HMC JIT-compile blocker by switching
to \texttt{PreconditionedNoUTurnSampler} with a static
\texttt{max\_tree\_depth} and no \texttt{pmap} wrapper, but the
\texttt{DualAveragingStepSizeAdaptation} adapter collapses on this
model: the local-curvature mismatch between modes prevents a single
mass matrix or step size from working, and a windowed mass-matrix
adapter would require a \texttt{JointDistribution} refactor of the
\gigalens\ prob model. The v11f NUTS run, started from the v7
paper-mode MAP point, converges into a {\it third}\ distinct local
mode at $\log p\approx -63{,}868$ with $\thetaE=2.604''$ ($1.6\%$ off
the published value, our closest reproduction of \thetaE), confirming
that the posterior surface is genuinely multimodal but not yielding a
mixed chain.
\item We document six upstream software issues in \gigalens\,2.0/
JAX\,0.6.2/TFP\,0.25.0 that would, if resolved, materially improve
the reproducibility of any Foundry-lineage analysis --- notably the
silent NaN from \texttt{jnp.linalg.cholesky} on a rank-deficient SVI
covariance, and the dual-averaging step-size pathology on
multimodal-Hessian models.
\end{enumerate}

For the broader DESI Strong Lens Foundry program, and for the
GIGA-Lens lineage at large, these findings argue that an
in-exposure empirical PSF and a multi-start MAP scan are reasonable
defaults to add to the published recipe, that mass--light PA tying
or an equivalent multipole correction is a worthwhile next-order
modeling improvement, that NUTS with static \texttt{max\_tree\_depth}
should replace HMC as the default JIT-tractable sampler at this model
size, and that mode-aware adaptation (windowed mass-matrix estimation,
basin-conditional re-tuning, or transparent reporting of the basin a
chain occupies) is a prerequisite for trustworthy posterior inference
at LSST/Roman/Euclid scale.

\begin{acknowledgments}
This work made use of \emph{Hubble Space Telescope}\ observations
collected under program GO--15867 (PI Huang) at the Space Telescope
Science Institute, which is operated by the Association of
Universities for Research in Astronomy, Inc., under NASA contract
NAS5--26555, and retrieved from MAST (DOI 10.17909/hx0v-9260). We
thank the \gigalens\ authors \citep{gu2022gigalens} for releasing
their differentiable forward-modeling framework as open source, and
the DESI Collaboration \citep{desicollab2016} for the larger lensing
program of which Foundry I is a part. Computational resources were
provided by the Department of Computer Science, University of
San Francisco. X.H.\ thanks Lawrence Berkeley National Laboratory for
ongoing support.
\end{acknowledgments}

\software{
\gigalens\ \citep{gu2022gigalens},
\lenstronomy\ \citep{birrer2018lenstronomy},
JAX \citep{bradbury2018jax},
TensorFlow Probability \citep{dillon2017tfp},
\photutils\ \citep{bradley2024photutils},
Astropy \citep{astropy2022}.
}

%%%%%%%%%%%%%%%%%%%%%%%%%%%%%%%%%%%%%%%%%%%%%%%%%%%%%%%%%%%%%%%%%%%%%%%%%%%%%%%
\appendix
%%%%%%%%%%%%%%%%%%%%%%%%%%%%%%%%%%%%%%%%%%%%%%%%%%%%%%%%%%%%%%%%%%%%%%%%%%%%%%%

\section{Internal Ablation: v1 $\rightarrow$ v11}
\label{app:ablation}

Table~\ref{tab:ablation} traces the per-version mass-fit accuracy across
the internal versions of this reproduction. The first four versions
(v1--v4) explore minimal-model variants without shapelets or with
unconstrained lens-light intensities. v5 introduces positive-$I_{e}$
\texttt{LogNormal} priors and a tight mass-center prior, and brings
\thetaE\ within $0.5\%$ of the paper. v6 adds the shapelet basis with
explicit per-coefficient priors and is the smallest model that produces
a physically reasonable global $\chi^{2}$. v7 is a 200-chain multi-start
MAP scan (TinyTim PSF) that exposes the multimodality of the surface.
v8 is an SVI posterior initialised from the v7 paper-mode chain, still
with TinyTim. v9 repeats v7's multi-start with the empirical PSF,
flipping the global best into paper-mode. v10 is the SVI posterior
initialised from the v9 global-best paper-mode chain and is the
canonical headline posterior of this paper (Table~\ref{tab:mass}).
v11 is a NUTS-based diagnostic series (\S\ref{sec:nuts}) that resolves
the HMC JIT-compile blocker, exposes a silent NaN in
\texttt{jnp.linalg.cholesky} on the rank-deficient SVI covariance,
demonstrates that \texttt{DualAveragingStepSizeAdaptation} collapses
on this model, and --- in the v11f run --- converges into a third local
mode with $\thetaE=2.604''$ ($1.6\%$ off the published value, our
closest single-version match for \thetaE).

\begin{deluxetable*}{lccccccc}
\tablecaption{Ablation summary v1$\to$v11. Best-fit mass-parameter
values for each version, paper values for reference.\label{tab:ablation}}
\tablehead{
\colhead{Version} & \colhead{\thetaE} & \colhead{\gammaEPL} & \colhead{\eone} &
\colhead{\etwo} & \colhead{\gextone} & \colhead{\gexttwo} & \colhead{$\chi^{2}_{\mathrm{best}}$}
}
\startdata
v1   & $+2.697$  & $+1.66$ & $-0.015$ & $+0.011$ & $+0.001$ & $+0.053$ & $24.65$ \\
v2   & $+2.062$  & $+1.57$ & $+0.141$ & $-0.108$ & $-0.001$ & $+0.032$ & $7.98$  \\
v4   & $+1.523$  & $+1.47$ & $-0.035$ & $-0.076$ & $-0.036$ & $-0.064$ & $3.45$ \tablenotemark{a}  \\
v5   & $+2.658$  & $+1.72$ & $-0.102$ & $-0.179$ & $-0.059$ & $-0.110$ & $7.56$  \\
v6   & $+2.632$  & $+1.73$ & $+0.128$ & $+0.073$ & $+0.065$ & $+0.005$ & $6.59$  \\
v8 (SVI/TT) & $+2.619$ & $+2.20$ & $+0.146$ & $-0.150$ & $+0.044$ & $-0.106$ & $7.72$ \tablenotemark{b} \\
v9 best (MAP/eP)  & $+2.578$ & $+2.16$ & $+0.103$ & $-0.045$ & $+0.034$ & $-0.067$ & $11.93$ \\
v10 (SVI/eP)\tablenotemark{c}  & $+2.566$ & $+2.15$ & $+0.106$ & $-0.053$ & $+0.040$ & $-0.068$ & --- \\
v11f (NUTS/eP)\tablenotemark{d} & $+2.604$ & $+2.26$ & $+0.159$ & $-0.151$ & $+0.054$ & $-0.119$ & --- \\
\hline
Paper (HMC)         & $+2.6463$ & $+1.372$ & $+0.1091$ & $-0.1320$ & $+0.0657$ & $-0.0939$ & --- \\
\enddata
\tablenotetext{a}{v4 has the lowest $\chi^{2}$ but is unphysical: the
unconstrained \texttt{use\_lstsq=True} setting allows negative
lens-light intensities; the resulting Einstein radius is $43\%$ low.}
\tablenotetext{b}{v8 is the first SVI posterior; initialized from the
best paper-mode chain of the v7 200-chain multi-start MAP with TinyTim
PSF.}
\tablenotetext{c}{v10 is the canonical posterior reported in
Table~\ref{tab:mass}; initialized from the v9 global-best paper-mode
chain (empirical PSF, $\chi^{2}=11.93$). v10 SVI does not produce a
new $\chi^{2}$ value; the entry is left blank.}
\tablenotetext{d}{v11f is the only NUTS variant of the v11 series
(\S\ref{sec:nuts}) that converged into a basin distinct from the
v10 SVI mode; it sits at $\log p\approx -63{,}868$ and reproduces
\thetaE\ to $1.6\%$ but the chain did not mix and is reported as a
second-mode point estimate, not a posterior. NUTS resolves the
HMC JIT-compile blocker (v11) but \texttt{DualAveragingStepSizeAdaptation}
collapses step size to $\sim$2--$3\times10^{-6}$ on this model
(v11a/e/f), so the chain barely moves once in a basin.}
\tablecomments{TT $=$ TinyTim PSF. eP $=$ empirical PSF. ``Paper-mode''
is the sign quadrant $(\eone{>}0,\etwo{<}0,\gextone{>}0,\gexttwo{<}0)$.
The flip in sign quadrant at v6$\to$v7 reflects multimodality, not
a coding regression: v6 is a single-chain MAP at the global best
under TinyTim, which lies in a different quadrant than the paper.
v11f confirms a third distinct basin reached by gradient-based
sampling.}
\end{deluxetable*}

\section{Prior Specification}
\label{app:priors}

Table~\ref{tab:priors} reproduces our prior distributions, which
closely follow \citet[][Table 2]{huang2025foundry1} with the
modifications discussed in \S\ref{sec:model}.

\begin{deluxetable}{lll}
\tablecaption{Prior specification on the 74 free non-linear parameters
of the lens model.\label{tab:priors}}
\tablehead{
\colhead{Parameter group/parameter} & \colhead{Distribution} &
\colhead{Parameters}
}
\startdata
\hline
\multicolumn{3}{l}{\textit{Mass (EPL + external shear; 8 params)}} \\
\thetaE\        & LogNormal       & $\mu=\log 2.5$, $\sigma=0.25$ \\
\gammaEPL\      & TruncatedNormal & $\mu=2.0$, $\sigma=0.25$, $[1.0,2.7]$ \\
$\eone^{\rm m}$       & Normal          & $\mu=0$, $\sigma=0.1$ \\
$\etwo^{\rm m}$       & Normal          & $\mu=0$, $\sigma=0.1$ \\
$c_{x}^{\rm mass}$ & Normal       & $\mu=0$, $\sigma=0.02$\,arcsec \\
$c_{y}^{\rm mass}$ & Normal       & $\mu=0$, $\sigma=0.02$\,arcsec \\
\gextone\       & Normal          & $\mu=0$, $\sigma=0.05$ \\
\gexttwo\       & Normal          & $\mu=0$, $\sigma=0.05$ \\
\hline
\multicolumn{3}{l}{\textit{Lens light --- main lens (2$\times$ S\'{e}rsic, 14 params)}} \\
$R_{\mathrm{eff}}$ (compact)  & LogNormal & $\mu=\log 0.4$, $\sigma=0.3$ \\
$R_{\mathrm{eff}}$ (extended) & LogNormal & $\mu=\log 2.0$, $\sigma=0.3$ \\
$n_{\rm S}$       & Uniform         & $[0.5, 8.0]$ \\
$e_{1,2}^{\rm light}$  & TruncatedNormal & $\mu=0$, $\sigma=0.15$, $[-0.4,0.4]$ \\
$c_{x},c_{y}$     & Normal          & $\mu=0$, $\sigma=0.02$\,arcsec \\
$I_{e}$ (compact) & LogNormal       & $\mu=\log 5.0$, $\sigma=0.5$ \\
$I_{e}$ (extended)& LogNormal       & $\mu=\log 2.0$, $\sigma=0.5$ \\
\hline
\multicolumn{3}{l}{\textit{Lens light --- nearby companion (2$\times$ S\'{e}rsic, 14 params)}} \\
$R_{\mathrm{eff}}$ & LogNormal & $\mu=\log 0.3,\,0.6$, $\sigma=0.3$ \\
$c_{x},c_{y}$ & Normal & $\mu=(-2.34,-2.86)$, $\sigma=0.1$\,arcsec \\
$I_{e}$ & LogNormal & $\mu=\log 1.0,\,0.5$, $\sigma=0.5$ \\
$e,n$ & As main lens & --- \\
\hline
\multicolumn{3}{l}{\textit{Source S\'{e}rsic (7 params)}} \\
$R_{\mathrm{eff}}$ & LogNormal & $\mu=\log 0.5$, $\sigma=0.3$ \\
$n_{\rm S}$ & Uniform & $[0.5, 6.0]$ \\
$e_{1,2}^{\rm src}$ & TruncatedNormal & $\mu=0$, $\sigma=0.15$, $[-0.5,0.5]$ \\
$c_{x},c_{y}$ & Normal & $\mu=0$, $\sigma=0.1$\,arcsec \\
$I_{e}$ & LogNormal & $\mu=\log 2.0$, $\sigma=0.5$ \\
\hline
\multicolumn{3}{l}{\textit{Source shapelets ($n_{\max}=6$, $1+2+28$ params)}} \\
$\beta$ & LogNormal & $\mu=\log 0.1$, $\sigma=0.1$ \\
$c_{x},c_{y}$ & Normal & $\mu=0$, $\sigma=0.05$\,arcsec \\
amp$_{i}$, $i=0,\dots,27$ & Normal & $\mu=0$, $\sigma=5/\sqrt{i+1}$ \\
\enddata
\tablecomments{The shapelet amplitude scale follows the gigalens
shapelets demonstration. The companion-light centroid prior centers
on the source-extraction detection on the v2 residual map
(Figure~\ref{fig:nearby}). Total: 74 free non-linear parameters.}
\end{deluxetable}

\bibliographystyle{aasjournalv7}
\bibliography{references}

\end{document}
